\subsubsection{Paralog Analysis}
The Python module provides comprehensive wrappers for paralog analysis.

\textbf{Key Features:}
\begin{itemize}
    \item \textbf{Detect Neofunctionalization}: Identifies neofunctionalization for genes by checking whether the difference of expression to its ancestor exceeds the threshold for the respective axis.
    \item \textbf{Detect Dosage Effect}: Identifies subsets of paralogs with small angle to the \lstinline!ancestor! (\lstinline!max_angle!) and sum to a magnitude significantly exceeding \lstinline!norm(ancestor)! (\lstinline!gain!).
    \item \textbf{Detect Subfunctionalization}: Identifies subsets of paralogs exhibiting significant angles to the \lstinline!ancestor!.
\end{itemize}

\begin{enumerate}
    \item \textbf{Detect Neofunctionalization}

        Identify neofunctionalization for genes by checking whether the difference of expression to its \lstinline!ancestor! exceeds the threshold for each axis.

        \textbf{Arguments:}
        \begin{itemize}
            \item \lstinline!ancestors (np.ndarray)!: 2D array of shape \lstinline!(n_axes, n_families)!.
            \item \lstinline!genes (np.ndarray)!: 2D array of shape \lstinline!(n_axes, n_genes)!.
            \item \lstinline!gene_to_fam (np.ndarray)!: 1D integer array of length \lstinline!n_genes! mapping gene index to family index.
            \item \lstinline!thresholds (np.ndarray)!: 1D array of shape \lstinline!(n_axes,)! with per-axis thresholds.
        \end{itemize}

        \textbf{Returns:}
        \begin{itemize}
            \item \lstinline!np.ndarray!: Integer array of shape \lstinline!(n_genes, n_axes)! with 0/1 values  
                  (non-zero interpreted as \lstinline!True!).
        \end{itemize}

\begin{lstlisting}[language=Python, caption={Neofunctionalization Detection}]
neofunc_mask = tox_detect_neofunctionalization(ancestors, genes, gene_to_fam, thresholds)
\end{lstlisting}


    \item \textbf{Detect Dosage Effect}

        Detect dosage effect among paralogs using gain and angle thresholds.

        This function has binomial memory complexity, depending on included genes. This is why this subroutine only handles one single family, the size varies too much across them. If it is wanted to have all results in one array, just use a python list.

        \textbf{Arguments:}
        \begin{itemize}
            \item \lstinline!ancestor (array-like)!: vector of shape \lstinline!(n_dims)!.
            \item \lstinline!genes (array-like)!: 2D array of shape \lstinline!(n_dims, n_genes)! containing all genes.
            \item \lstinline!filtered_paralogs_mask (array-like)!: int32 mask (chunks of 32 bits) for \lstinline!n_genes!.
            \item \lstinline!max_subset_size (int)!: desired maximum subset size (may be adjusted).
            \item \lstinline!gain_gamma (float)!: required magnitude gain (default \lstinline!0.1!).
            \item \lstinline!max_angle (float)!: maximum allowed angle in radians (default \(\pi\)).
        \end{itemize}

        \textbf{Returns:}
        \begin{itemize}
            \item \lstinline!dict!: \{ \lstinline!'n_results': int!, \lstinline!'results': np.ndarray! \}
        \end{itemize}

\begin{lstlisting}[language=Python, caption={Dosage Effect Detection}]
n_results, results = tox_detect_dosage_effect(
    ancestor, genes, filtered_paralogs_mask, max_subset_size,
    gain_gamma=0.1, max_angle=np.pi
).values()
\end{lstlisting}


    \item \textbf{Detect Subfunctionalization}

        Detect subfunctionalization among paralogs based on residual distance and pruning.

        This function has binomial memory complexity, depending on included genes. This is why this subroutine only handles one single family, the size varies too much across them. If it is wanted to have all results in one array, just use a python list.

        \textbf{Arguments:}
        \begin{itemize}
            \item \lstinline!ancestor (array-like)!: vector of shape \lstinline!(n_dims)!.
            \item \lstinline!genes (array-like)!: 2D array of shape \lstinline!(n_dims, n_genes)! containing all genes.
            \item \lstinline!rdi_threshold (float)!: maximum allowed residual distance to \lstinline!ancestor!.
            \item \lstinline!filtered_paralogs_mask (array-like)!: int32 mask (chunks of 32 bits) for \lstinline!n_genes!.
            \item \lstinline!max_subset_size (int)!: desired maximum subset size (may be adjusted).
            \item \lstinline!paralog_norms (array-like)!: float64 norms for all genes (\lstinline!n_genes!).
            \item \lstinline!sorted_paralog_norms_perm (array-like)!: int32 permutation of norms (\lstinline!n_genes!).
        \end{itemize}

        \textbf{Returns:}
        \begin{itemize}
            \item \lstinline!dict!: \{ \lstinline!'n_results': int!, \lstinline!'results': np.ndarray! \}
        \end{itemize}

\begin{lstlisting}[language=Python, caption={Subfunctionalization Detection}]
n_results, results = tox_detect_subfunctionalization(
    ancestor, genes, rdi_threshold,
    filtered_paralogs_mask, max_subset_size,
    paralog_norms, sorted_paralog_norms_perm
).values()
\end{lstlisting}


    \item \textbf{Mask Chunk Count}

        Compute the number of 32-bit chunks needed to encode a given number of genes.

        \textbf{Arguments:}
        \begin{itemize}
            \item \lstinline!n_genes (int)!: number of genes to encode.
        \end{itemize}

        \textbf{Returns:}
        \begin{itemize}
            \item \lstinline!int!: number of 32-bit chunks required.
        \end{itemize}

\begin{lstlisting}[language=Python, caption={Mask Chunk Count Calculation}]
n_chunks = tox_mask_chunk_count(n_genes)
\end{lstlisting}


    \item \textbf{Filter Paralogs By Pattern Subfunctionalization}

        Filter paralogs by subfunctionalization pattern using angle threshold, within the grouped slice.

        \textbf{Arguments:}
        \begin{itemize}
            \item \lstinline!gene_angles (array-like)!: angles for all genes (length = \lstinline!n_genes!).
            \item \lstinline!threshold (float)!: filtering threshold.
            \item \lstinline!gene_to_fam (array-like)!: gene-to-family mapping (\lstinline!n_genes!).
            \item \lstinline!n_families (int)!: number of families.
        \end{itemize}

        \textbf{Returns:}
        \begin{itemize}
            \item \lstinline!np.ndarray!: int32 mask (chunks of 32 bits) with 1 for kept, 0 otherwise.
        \end{itemize}

\begin{lstlisting}[language=Python, caption={Subfunctionalization Pattern Filtering}]
mask = tox_filter_paralogs_by_pattern_subfunctionalization(
    gene_angles, threshold, gene_to_fam, n_families
)
\end{lstlisting}



    \item \textbf{Filter Paralogs By Pattern Dosage Effect}

        Filter paralogs by dosage effect using angle threshold, within the grouped slice.

        \textbf{Arguments:}
        \begin{itemize}
            \item \lstinline!gene_angles (array-like)!: angles for all genes (length = \lstinline!n_genes!).
            \item \lstinline!threshold (float)!: filtering threshold.
            \item \lstinline!gene_to_fam (array-like)!: gene-to-family mapping (\lstinline!n_genes!).
            \item \lstinline!n_families (int)!: number of families.
        \end{itemize}

        \textbf{Returns:}
        \begin{itemize}
            \item \lstinline!np.ndarray!: int32 mask (chunks of 32 bits) with 1 for kept, 0 otherwise.
        \end{itemize}

\begin{lstlisting}[language=Python, caption={Dosage Effect Pattern Filtering}]
mask = tox_filter_paralogs_by_pattern_dosage_effect(
    gene_angles, threshold, gene_to_fam, n_families
)
\end{lstlisting}


    \item \textbf{Calculate Work Array Paralog Subsets Size}

        Calculate the required work array size for paralog subset analysis.

        \textbf{Arguments:}
        \begin{itemize}
            \item \lstinline!max_subset_size (int)!: maximum subset size (may be adjusted by the routine).
            \item \lstinline!n_genes (int)!: total number of genes.
            \item \lstinline!filtered_paralogs_mask (array-like)!: integer bit mask (chunks of 32 bits) marking filtered paralogs.
        \end{itemize}

        \textbf{Returns:}
        \begin{itemize}
            \item \lstinline!dict!: \{ \lstinline!'actual_max_subset_size': int!, \lstinline!'work_array_size': int! \}
        \end{itemize}

\begin{lstlisting}[language=Python, caption={Work Array Size Calculation}]
actual_max_subset_size, work_array_size = tox_calc_work_arr_paralog_subsets_size(
    max_subset_size, n_genes, filtered_paralogs_mask
).values()
\end{lstlisting}


    \item \textbf{Mask Check State}

        Compute the number of 32-bit chunks needed to encode a given number of genes.

        \textbf{Arguments:}
        \begin{itemize}
            \item \lstinline!n_genes (int)!: number of genes to encode.
        \end{itemize}

        \textbf{Returns:}
        \begin{itemize}
            \item \lstinline!int!: number of 32-bit chunks required.
        \end{itemize}

\begin{lstlisting}[language=Python, caption={Mask Chunk Count Calculation}]
n_chunks = tox_mask_chunk_count(n_genes)
\end{lstlisting}


\end{enumerate}